% Generated draft: requires mathematical and source review.
\documentclass[11pt]{article}
\usepackage[margin=25mm]{geometry}
\usepackage{fontspec}
\setmainfont{Latin Modern Roman}
\usepackage{amsmath,amssymb,mathtools,unicode-math}
\setmathfont{Latin Modern Math}
\usepackage{enumitem,needspace,xurl,hyperref}
\setlength{\parindent}{0pt}
\setlength{\parskip}{5pt}
\setlength{\emergencystretch}{3em}
\newcommand{\N}{\mathbb{N}}
\newcommand{\Z}{\mathbb{Z}}
\newcommand{\Q}{\mathbb{Q}}
\newcommand{\R}{\mathbb{R}}
\newcommand{\C}{\mathbb{C}}
\newcommand{\F}{\mathbb{F}}
\newcommand{\A}{\mathbb{A}}
\newcommand{\PP}{\mathbb{P}}
\newcommand{\E}{\mathbb{E}}
\newcommand{\eps}{\varepsilon}
\newcommand{\dd}{\,\mathrm{d}}
\begin{document}
\section*{UnsolvedMath ID: 30000781}
\textbf{Source problem code:} catalogue entry
\subsection{Definitions and mathematical statement}
Fix a CM Hida deformation in the setting of the source report. Let $p$ be its prime and $K$ its CM field. The intended analytic $p$-adic $L$-functions are denoted
\[
 L_p^{\mathrm{an}}(\operatorname{GL}(1)/K),
 \qquad
 L_p^{\mathrm{an}}(\operatorname{GL}(2)/\Q).
\]
The conjecture asks whether their generated ideals coincide:
\[
 \bigl(L_p^{\mathrm{an}}(\operatorname{GL}(1)/K)\bigr)
   =\bigl(L_p^{\mathrm{an}}(\operatorname{GL}(2)/\Q)\bigr).
\]
This equality must be interpreted in the report's common coefficient/Iwasawa ring after its prescribed identifications or comparison maps. Neither that ring nor those maps are supplied, so equality of ideals from potentially different initial rings is not asserted without qualification.

The usual context is an ordinary $p$-adic Hida family with complex multiplication, related to a family of Hecke characters. The field $K$, characters, deformation, ordinarity and prime restrictions, interpolation variables, periods, Euler factors, and analytic constructions are not defined in the extraction. No imaginary-quadratic, splitting, residual, or normalization hypothesis is silently selected solely from the notation. The mathematical request is equality of ideals, not literal equality of normalized functions. In a common integral domain with nonzero generators, equality would mean equality up to a unit; this explanation is conditional on the missing ring hypotheses.

The original states that this is equivalent to the corresponding CM Iwasawa main conjecture. That equivalence is retained as a supplied assertion, but no Selmer object, characteristic ideal, algebraic $L$-function, or hypotheses supporting the equivalence are provided. The assessment reports a proof under irreducibility and vanishing of a $\mu$-invariant, without identifying which representation or invariant. Those conditions describe a supplied partial case and are not inserted into the unrestricted conjectural question.
\subsection{Short English statement}
The question compares the principal ideals generated by the two analytic $p$-adic $L$-function constructions for a CM Hida deformation. The common ring, comparison maps, normalizations, and the claimed main-conjecture equivalence are missing; the reported irreducibility and $\mu=0$ case remains attributed partial context.
\subsection{Sources}
\begin{itemize}
\item Local UnsolvedMath registry, permanent ID 30000781, code catalogue entry, ``CM Hida Families and $p$-Adic $L$-Function Ideals.'' Proposed year 2007; proposer unspecified.
\item \emph{Algebraische Zahlentheorie} (2008), Oberwolfach Reports, pp.\ 1741--1792, DOI 10.4171/owr/2007/30; \url{https://doi.org/10.4171/owr/2007/30}. The displayed report year 2008 differs from the DOI and proposed year 2007. This is the sole supplied literature reference.
\item The original additionally asserts equivalence with the corresponding CM Iwasawa main conjecture, without supplying its formulation.
\item Statement status corrected\_verified; the supplied note says the analytic-ideal equality is explicit and notation normalized.
\item Literature triage checked 2026-08-21: partially solved. It reports the comparison proved under irreducibility and vanishing of a $\mu$-invariant, leaving the unrestricted question unresolved. No exact representation, invariant, theorem, or proof citation is given beyond the report. Created 2026-08-25; updated 2026-09-16. Current status and precise hypotheses require independent verification.
\item Registry curation by Ulam AI / UnsolvedMath Contributors, CC BY 4.0; the underlying report retains its supplied provenance. Unreviewed statement-only draft.
\end{itemize}
\end{document}
